\chapter{APDU over T=0}
\label{app:apdu-over-t0}

This appendix summarizes how the present implementation maps the
conceptual APDU model of Chapter~4 onto the concrete \texttt{T=0}
transport and onto the programming interface exposed by Oxide SE.

From an implementation point of view, this mapping lives in a small
stack of kernel objects:
\begin{itemize}
\item \texttt{TransportLayer} owns byte I/O;
\item \texttt{T0ApduManager<T>} owns the \texttt{T=0} state machine;
\item \texttt{ApduLayer} owns protocol-layer composition above
  \texttt{T=0};
\item \texttt{SEApdu} remains the card-side command-processing
  interface.
\end{itemize}

This means that the appendix should be read at two levels at once:
\begin{itemize}
\item as a description of the external \texttt{T=0} byte stream;
\item as a description of how that byte stream is reified into layered
  kernel objects before command logic runs.
\end{itemize}

\section{From Conceptual APDUs to the Oxide SE Interface}

At the conceptual level, an APDU is a command/response pair made of:
\begin{itemize}
\item one command header (\texttt{CLA}, \texttt{INS}, \texttt{P1},
  \texttt{P2});
\item optional incoming command data, logically measured by
  \texttt{Lc};
\item optional outgoing response data, logically measured by
  \texttt{Le};
\item one final status word \texttt{SW1/SW2}.
\end{itemize}

At the implementation level, Oxide SE does not expose that conceptual
model as a ready-made ``full APDU'' object that already contains both
incoming and outgoing data. Instead, it exposes the APDU in the form in
which a secure element can actually process it under \texttt{T=0}:
\begin{itemize}
\item first, the five command bytes are known;
\item then the command logic decides whether command data must be
  received;
\item then the command logic decides whether response data will be
  produced, and how many bytes are logically available.
\end{itemize}

This is why the common APDU-facing interface is organized around the
following operations:
\begin{itemize}
\item \texttt{set\_incoming\_and\_Receive()};
\item \texttt{set\_outgoing()};
\item \texttt{set\_outgoing\_length(len)};
\item \texttt{send\_bytes(offset, len)};
\item \texttt{send\_bytes\_long(src, offset, len)}.
\end{itemize}

These operations are not transport bypasses. They are declarations and
state transitions on the card-side APDU session:
\begin{itemize}
\item \texttt{set\_incoming\_and\_Receive()} means ``this command consumes
  incoming data now'';
\item \texttt{set\_outgoing()} means ``this command is logically an
  outgoing command'';
\item \texttt{set\_outgoing\_length(len)} means ``the command has prepared
  \texttt{len} outgoing bytes'';
\item \texttt{send\_bytes()} and \texttt{send\_bytes\_long()} only stage
  payload bytes in the APDU buffer and set the outgoing length
  accordingly.
\end{itemize}

The final transport behavior is still chosen by the kernel-owned
\texttt{T=0} loop after the command handler returns.

\section{From the Interface to the \texttt{T=0} Exchange}

The implementation distinguishes the four classical short APDU cases:
\begin{itemize}
\item \textbf{Case 1:} no incoming data, no outgoing data;
\item \textbf{Case 2:} no incoming data, outgoing data only;
\item \textbf{Case 3:} incoming data only;
\item \textbf{Case 4:} incoming data followed by a deferred outgoing
  response.
\end{itemize}

The present interface maps those cases as follows:

\begin{center}
\begin{tabular}{|p{2.1cm}|p{6.0cm}|p{6.0cm}|}
\hline
\textbf{Case} & \textbf{Oxide SE operations} & \textbf{Transport result} \\
\hline
Case 1 & no APDU transition call & final \texttt{SW1/SW2} \\
\hline
Case 2 & \texttt{set\_outgoing()}\newline
         \texttt{set\_outgoing\_length()} &
         procedure byte, data, final status \\
\hline
Case 3 & \texttt{set\_incoming\_and\_Receive()} &
         procedure byte, incoming bytes, final status \\
\hline
Case 4 & \texttt{set\_incoming\_and\_Receive()}\newline
         \texttt{set\_outgoing()}\newline
         \texttt{set\_outgoing\_length()} &
         incoming phase, then \texttt{61xx}, then \texttt{GET RESPONSE} \\
\hline
\end{tabular}
\end{center}

This is the main implementation rule of the current system: APDU command
logic expresses \emph{what kind of APDU interaction is required}, while
the transport loop remains responsible for deciding \emph{which bytes
must actually be emitted next on the wire}.

In implementation terms:
\begin{itemize}
\item \texttt{SEApdu} expresses the card-side command intent;
\item \texttt{ApduLayer} transports command and completion objects
  between protocol layers;
\item \texttt{T0ApduManager} translates those decisions into concrete
  \texttt{T=0} bytes.
\end{itemize}

\section{The Current \texttt{T=0} State Machine}

The current implementation can be summarized by the following state
machine:

\begin{enumerate}
\item emit ATR after reset;
\item read the 5-byte command header;
\item emit one \texttt{0x60} NULL procedure byte to restart the work
  waiting time without yet committing to an incoming or outgoing data
  phase, as allowed by the \texttt{T=0} procedure-byte rules of
  ISO/IEC~7816-3 (\emph{Protocol type T=0, structure and processing of
  commands}, clause \texttt{3.5.b}) \cite{iso7816-3};
\item build the transport-owned APDU object;
\item dispatch the command logic;
\item depending on the declared APDU direction:
  \begin{itemize}
  \item return only the final status;
  \item emit a procedure byte, then direct outgoing bytes, then final
    status;
  \item emit \texttt{61xx} and store the pending response;
  \item emit \texttt{6Cxx} and wait for command replay with the correct
    \texttt{Le}.
  \end{itemize}
\item if a valid \texttt{GET RESPONSE} command is later received, emit
  the \texttt{C0} procedure byte, the next response chunk, and either a
  final status or another \texttt{61xx}.
\end{enumerate}

The critical point is that the command handler does not itself drive the
wire. It only mutates APDU state. The kernel transport loop interprets
that state and performs the corresponding \texttt{T=0} exchange.

The current reference kernel instantiates this state machine through the
following object composition:
\begin{center}
\texttt{CurrentTransport $\rightarrow$ T0ApduManager $\rightarrow$ PassthroughApduLayer $\rightarrow$ SecureChannelLayer $\rightarrow$ TracingApduLayer}
\end{center}

The dispatcher then consumes the resulting \texttt{ApduCommand},
derives an \texttt{SEApdu} view from it, and finally reports an
\texttt{ApduCompletion} back into the same stack.

\section{Procedure Bytes}

Procedure bytes are the bridge between the conceptual APDU model and the
half-duplex \texttt{T=0} byte stream.

In the present implementation:
\begin{itemize}
\item immediately after the 5-byte header, the secure element emits the
  NULL procedure byte \texttt{0x60} once, in order to restart the work
  waiting time while command routing begins; this is the dedicated
  ``wait'' procedure byte defined by ISO/IEC~7816-3
  clause~\texttt{3.5.b} and it requests no data transfer by itself
  \cite{iso7816-3};
\item for incoming data reception, the secure element emits the current
  command \texttt{INS} as a procedure byte before receiving the
  \texttt{Lc} bytes;
\item for direct outgoing transmission, the secure element emits the
  current command \texttt{INS} as a procedure byte before sending the
  response bytes;
\item for \texttt{GET RESPONSE}, the secure element emits
  \texttt{C0}---the \texttt{INS} byte of \texttt{GET RESPONSE}---before
  sending the deferred response bytes.
\end{itemize}

This means that the APDU programming interface intentionally does not
contain an explicit ``send procedure byte'' operation. Procedure bytes
are not an application-level concept in Oxide SE. They are an effect of
the kernel transport state machine. In particular, Oxide SE distinguishes:
\begin{itemize}
\item \texttt{0x60} as a NULL procedure byte used only to extend the
  waiting time;
\item \texttt{INS} or \texttt{C0} as whole-transfer ACK procedure bytes used only when
  the transport actually enters an incoming, direct outgoing, or
  \texttt{GET RESPONSE} phase.
\end{itemize}

The protocol also defines the one's complement of \texttt{INS} as a
single-byte-transfer ACK. The reference kernel currently emits only the
whole-transfer ACK, but host parsers distinguish the complemented ACK
from both status bytes and transport errors. Conversely, receipt of a
valid \texttt{SW1} immediately after the header terminates the command
before any incoming or outgoing data phase. The legal byte classes and
the resulting \texttt{INS} restriction are specified in
Section~\ref{sec:apdu-t0-procedure-bytes}.

\section{Wrong Length and Deferred Response}

Two transport-owned completion mechanisms are particularly important.

\subsection{\texttt{6Cxx} for Pure Outgoing Commands}

For a pure outgoing command, the fifth command byte is interpreted as
\texttt{Le}. The command logic may declare an actual response length
through \texttt{set\_outgoing\_length(len)}. If the requested
\texttt{Le} does not match the actual available length, the current
implementation may return \texttt{6Cxx}, where \texttt{xx} is the
correct length to request.

The host must then resend the same command with the corrected
\texttt{Le}. This path is therefore transport-owned: the command logic
does not manually emit \texttt{6Cxx}.

\subsection{\texttt{61xx} and \texttt{GET RESPONSE} for In/Out Commands}

For a command that first consumes incoming data and then prepares
outgoing data, the current implementation follows the classical
\texttt{T=0} pattern:
\begin{itemize}
\item the command consumes incoming data through
  \texttt{set\_incoming\_and\_Receive()};
\item the command prepares outgoing data through
  \texttt{set\_outgoing()} and \texttt{set\_outgoing\_length(len)};
\item the kernel does not emit those outgoing bytes immediately;
\item instead, it stores them in the deferred-response buffer and
  returns \texttt{61xx};
\item the interface device must then send \texttt{GET RESPONSE};
\item the kernel emits the \texttt{C0} procedure byte, then the stored
  response bytes, then the final \texttt{SW1/SW2}.
\end{itemize}

This is the current realization of a logical Case~4 APDU in Oxide SE.

\section{Synthesis}

The implementation can therefore be summarized in one sentence:
Oxide SE exposes APDUs to command logic as a Java-Card-like card-side
interface, while the kernel remains the sole owner of the concrete
\texttt{T=0} byte stream, including procedure bytes, \texttt{6Cxx},
\texttt{61xx}, and \texttt{GET RESPONSE}.

This separation allows the same command-processing model to be used:
\begin{itemize}
\item by kernel-side code through the kernel \texttt{Apdu} view;
\item by Rustlets through \texttt{RustletCtx};
\item by the kernel-to-Rustlet bridge through the common
  \texttt{SEApdu} trait.
\end{itemize}
